\documentclass[a4paper,14pt]{extarticle}
\usepackage[utf8]{inputenc}
\usepackage[T2A]{fontenc}
\usepackage[russian]{babel}
\usepackage{amsmath,amssymb,amsfonts}
\usepackage{graphicx}
\usepackage{booktabs}
\usepackage{longtable}
\usepackage{listings}
\usepackage{xcolor}
\usepackage{hyperref}
\usepackage{geometry}
\usepackage{indentfirst}
\usepackage{tikz}
\usepackage{enumitem}

\geometry{left=3cm,right=1.5cm,top=2cm,bottom=2cm}
\setlength{\parindent}{1.25cm}
\linespread{1.5}

\hypersetup{
    colorlinks=true,
    linkcolor=blue,
    citecolor=blue,
    urlcolor=blue
}

\lstset{
    basicstyle=\ttfamily\small,
    breaklines=true,
    frame=single,
    backgroundcolor=\color{gray!10},
    keywordstyle=\color{blue},
    commentstyle=\color{green!60!black},
    stringstyle=\color{red!80!black},
    numbers=left,
    numberstyle=\tiny\color{gray},
    tabsize=2
}

\begin{document}

\section{Архитектура KubeBao: техническая реализация}

\subsection{Обзор системы}

KubeBao~--- Kubernetes-оператор для управления секретами через OpenBao (форк HashiCorp Vault). Система обеспечивает:

\begin{itemize}
    \item \textbf{Шифрование секретов etcd} через KMS Plugin API v2 (Kubernetes 1.25+) с использованием блочного шифра <<Кузнечик>> (ГОСТ~Р~34.12-2015);
    \item \textbf{Синхронизацию секретов} из OpenBao KV в Kubernetes Secrets через CRD \texttt{BaoSecret};
    \item \textbf{Декларативное управление политиками} OpenBao через CRD \texttt{BaoPolicy};
    \item \textbf{Монтирование секретов в поды} через Secrets Store CSI Driver.
\end{itemize}

Все криптографические операции выполняются собственной реализацией алгоритмов ГОСТ, без использования сторонних криптографических библиотек.

\subsection{Компонентная архитектура}

\begin{table}[h]
\centering
\caption{Компоненты системы KubeBao}
\begin{tabular}{@{}llp{8cm}@{}}
\toprule
\textbf{Компонент} & \textbf{Тип} & \textbf{Назначение} \\
\midrule
\texttt{kubebao-kms} & DaemonSet & KMS Plugin v2~--- шифрование/дешифрование секретов etcd \\
\texttt{kubebao-operator} & Deployment & Kubernetes Operator~--- управление CRD (BaoSecret, BaoPolicy) \\
\texttt{kubebao-csi} & DaemonSet & Secrets Store CSI Provider~--- монтирование секретов в файловую систему подов \\
\texttt{kubebao-ui} & Deployment & Веб-интерфейс управления и мониторинга \\
\bottomrule
\end{tabular}
\end{table}

\subsection{Математические основы криптографической подсистемы}

\subsubsection{Блочный шифр <<Кузнечик>> (ГОСТ~Р~34.12-2015)}

\paragraph{Параметры шифра.}

Блочный шифр <<Кузнечик>> определён стандартом ГОСТ~Р~34.12-2015 (также описан в RFC~7801). Основные параметры:

\begin{table}[h]
\centering
\begin{tabular}{@{}ll@{}}
\toprule
Параметр & Значение \\
\midrule
Размер блока ($n$) & 128 бит (16 байт) \\
Размер ключа ($k$) & 256 бит (32 байт) \\
Число раундов & 10 \\
Структура & SP-сеть (подстановочно-перестановочная) \\
\bottomrule
\end{tabular}
\end{table}

\paragraph{Алгебраическая структура.}

Шифр работает в поле Галуа $\mathrm{GF}(2^8)$ с неприводимым полиномом:
\begin{equation}
    p(x) = x^8 + x^7 + x^6 + x + 1
\end{equation}

Двоичное представление: \texttt{0x1C3}. Элементы поля~--- байты, операции (сложение~--- XOR, умножение~--- по модулю $p(x)$) используются в линейном преобразовании~$L$.

\paragraph{Нелинейное преобразование $S$ (подстановка).}

Нелинейное преобразование $S$ применяет фиксированную подстановку $\pi$ к каждому байту блока:
\begin{equation}
    S: V_{128} \to V_{128}, \quad S(a) = (\pi(a_{15}), \pi(a_{14}), \ldots, \pi(a_0))
\end{equation}
где $\pi: V_8 \to V_8$~--- фиксированная S-box таблица размером 256 элементов, определённая стандартом.

Свойства таблицы $\pi$:
\begin{itemize}
    \item высокая нелинейность (расстояние от аффинных функций $\geq 104$);
    \item алгебраическая степень 7;
    \item дифференциальная вероятность $\leq 2^{-6}$.
\end{itemize}

\paragraph{Линейное преобразование $L$.}

Линейное преобразование $L$ определяется через регистр сдвига с линейной обратной связью (LFSR) на основе функции~$l$:
\begin{equation}
    l(a) = \bigoplus_{i=0}^{15} k_i \cdot a_i \pmod{p(x)}
\end{equation}
где коэффициенты $k_i$~--- фиксированные элементы $\mathrm{GF}(2^8)$:

\begin{table}[h]
\centering
\small
\begin{tabular}{@{}c|cccccccccccccccc@{}}
\toprule
$i$ & 0 & 1 & 2 & 3 & 4 & 5 & 6 & 7 & 8 & 9 & 10 & 11 & 12 & 13 & 14 & 15 \\
\midrule
$k_i$ & 148 & 32 & 133 & 16 & 194 & 192 & 1 & 251 & 1 & 192 & 194 & 16 & 133 & 32 & 148 & 1 \\
\bottomrule
\end{tabular}
\end{table}

Преобразование $R$:
\begin{equation}
    R(a_{15}, \ldots, a_0) = (l(a_{15}, \ldots, a_0), a_{15}, a_{14}, \ldots, a_1)
\end{equation}

Преобразование $L$~--- 16-кратное применение $R$:
\begin{equation}
    L = R^{16}
\end{equation}

Это гарантирует MDS-свойство (Maximum Distance Separable): минимальное число активных S-box в любой дифференциальной характеристике составляет~17.

\paragraph{Развёртка ключа.}

Из 256-битного мастер-ключа $K = (K_1 \| K_2)$ вырабатываются 10 раундовых ключей $k_1, \ldots, k_{10}$:
\begin{equation}
    k_1 = K_1, \quad k_2 = K_2
\end{equation}

Для $i = 1, \ldots, 4$ (по 8 итераций Фейстеля на каждую пару):
\begin{equation}
    F[C_j](A, B) = (LSX[C_j](A) \oplus B, \; A)
\end{equation}
где $C_j$~--- итерационные константы, предвычисленные из $L$-преобразования:
\begin{equation}
    C_j = L(\mathrm{enc}(j))
\end{equation}

Функция $F$ реализует раунд сети Фейстеля:
\begin{equation}
    (A', B') = F[C](A, B) = (L(S(A \oplus C)) \oplus B, \; A)
\end{equation}

Каждая группа из 8 Фейстель-раундов порождает 2 новых раундовых ключа, итого 8 дополнительных ключей ($4 \text{ группы} \times 2 \text{ ключа}$).

\paragraph{Шифрование блока.}

Шифрование блока открытого текста $a \in V_{128}$:
\begin{align}
    a_0 &= a \\
    a_i &= LSX[k_i](a_{i-1}) = L(S(a_{i-1} \oplus k_i)), \quad i = 1, \ldots, 9 \\
    b &= X[k_{10}](a_9) = a_9 \oplus k_{10}
\end{align}
где:
\begin{itemize}
    \item $X[k](a) = a \oplus k$~--- сложение с раундовым ключом (Key Addition);
    \item $S$~--- нелинейная подстановка;
    \item $L$~--- линейное преобразование.
\end{itemize}

\paragraph{Табличная оптимизация реализации.}

В реализации KubeBao используется табличная оптимизация: совмещённые таблицы подстановки и линейного преобразования (Combined S-box + L-transform lookup). Для каждой из 16 байтовых позиций предвычислена таблица размером $256 \times 2$ (\texttt{uint64}):
\begin{equation}
    t[\mathrm{pos}][\mathrm{byte\_value}] = L\text{-transform contribution for } S(\mathrm{byte\_value}) \text{ at position pos}
\end{equation}

Это позволяет вычислить $L(S(\cdot))$ одним проходом XOR из 16 табличных значений, заменяя $16 \times 16$ операций умножения в поле Галуа на 16 обращений к таблице.

\subsubsection{Режим гаммирования CTR (ГОСТ~Р~34.13-2015, раздел~5.5)}

\paragraph{Определение.}

Режим CTR (Counter) превращает блочный шифр в потоковый. Согласно ГОСТ~Р~34.13-2015:

Зашифрование:
\begin{equation}
    C_i = P_i \oplus E_K(\mathrm{CTR}_i), \quad i = 1, \ldots, q
\end{equation}

Расшифрование:
\begin{equation}
    P_i = C_i \oplus E_K(\mathrm{CTR}_i), \quad i = 1, \ldots, q
\end{equation}
где:
\begin{itemize}
    \item $E_K$~--- функция зашифрования блочным шифром <<Кузнечик>> с ключом $K$;
    \item $P_i$~--- $i$-й блок открытого текста;
    \item $C_i$~--- $i$-й блок шифротекста;
    \item $\mathrm{CTR}_i$~--- $i$-й счётчик;
    \item $q = \lceil |P| / n \rceil$~--- число полных и неполных блоков.
\end{itemize}

\paragraph{Структура счётчика.}

128-битный счётчик разбивается на две половины:
\begin{equation}
    \mathrm{CTR} = \mathrm{CTR}_{\mathrm{upper}} \| \mathrm{CTR}_{\mathrm{lower}}
\end{equation}

\begin{itemize}
    \item \textbf{Верхние 64 бита} ($\mathrm{CTR}_{\mathrm{upper}}$)~--- инициализационный вектор (IV), фиксирован в пределах одного шифрования;
    \item \textbf{Нижние 64 бита} ($\mathrm{CTR}_{\mathrm{lower}}$)~--- инкрементируемый счётчик.
\end{itemize}

Функция инкремента (по ГОСТ~Р~34.13-2015, $s = n/2 = 64$):
\begin{equation}
    \mathrm{CTR}_{i+1} = \mathrm{CTR}_{\mathrm{upper}} \| (\mathrm{CTR}_{\mathrm{lower}} + 1 \bmod 2^{64})
\end{equation}

Верхние 64 бита \textbf{не изменяются} при инкременте~--- это гарантирует, что счётчик не может переполниться в пределах $2^{64}$ блоков ($\approx 295$ эксабайт для 128-битного блока).

\paragraph{Обработка последнего блока.}

Если длина открытого текста не кратна размеру блока, последний блок обрабатывается усечённой гаммой:
\begin{equation}
    C_q = P_q \oplus \mathrm{MSB}_{|P_q|}(E_K(\mathrm{CTR}_q))
\end{equation}
где $\mathrm{MSB}_r(\cdot)$~--- старшие $r$ бит.

\subsubsection{Алгоритм выработки имитовставки CMAC (ГОСТ~Р~34.13-2015, раздел~5.6)}

\paragraph{Определение.}

CMAC (Cipher-based Message Authentication Code) обеспечивает целостность и аутентичность данных. Определён в ГОСТ~Р~34.13-2015.

Длина выходного значения (тега): $|T| = n = 128$ бит.

\paragraph{Генерация подключей.}

Из ключа $K$ генерируются два подключа $K_1, K_2$:

\begin{enumerate}
    \item Вычислить $L = E_K(0^n)$~--- зашифрование нулевого блока.
    \item Подключ $K_1$:
    \begin{equation}
        K_1 = \begin{cases}
            L \ll 1 & \text{если } \mathrm{MSB}(L) = 0 \\
            (L \ll 1) \oplus R_b & \text{если } \mathrm{MSB}(L) = 1
        \end{cases}
    \end{equation}
    \item Подключ $K_2$:
    \begin{equation}
        K_2 = \begin{cases}
            K_1 \ll 1 & \text{если } \mathrm{MSB}(K_1) = 0 \\
            (K_1 \ll 1) \oplus R_b & \text{если } \mathrm{MSB}(K_1) = 1
        \end{cases}
    \end{equation}
\end{enumerate}
где:
\begin{itemize}
    \item $\ll 1$~--- побитовый сдвиг влево на 1 позицию;
    \item $R_b = 0^{120} \| 10000111_2 = \texttt{0x00\ldots0087}$~--- полином приведения для $\mathrm{GF}(2^{128})$.
\end{itemize}

Полином $R_b$ соответствует:
\begin{equation}
    x^{128} + x^7 + x^2 + x + 1
\end{equation}

Это минимальный неприводимый полином степени~128 над $\mathrm{GF}(2)$, что гарантирует максимальный период генератора подключей.

\paragraph{Вычисление CMAC.}

Для сообщения $M = M_1 \| M_2 \| \ldots \| M_q$:

\begin{enumerate}
    \item \textbf{Инициализация:} $X_0 = 0^n$
    \item \textbf{Обработка полных блоков} ($i = 1, \ldots, q-1$):
    \begin{equation}
        X_i = E_K(X_{i-1} \oplus M_i)
    \end{equation}
    \item \textbf{Обработка последнего блока:}
    \begin{itemize}
        \item Если последний блок полный ($|M_q| = n$):
        \begin{equation}
            T = E_K(X_{q-1} \oplus M_q \oplus K_1)
        \end{equation}
        \item Если последний блок неполный ($|M_q| < n$):
        \begin{equation}
            M_q^* = M_q \| 1 \| 0^{n - |M_q| - 1} \quad \text{(padding 10*)}
        \end{equation}
        \begin{equation}
            T = E_K(X_{q-1} \oplus M_q^* \oplus K_2)
        \end{equation}
    \end{itemize}
\end{enumerate}

В KubeBao CMAC вычисляется от конкатенации IV и шифротекста:
\begin{equation}
    T = \mathrm{CMAC}_{K_{\mathrm{mac}}}(\mathrm{IV} \| C)
\end{equation}

Это обеспечивает аутентификацию как шифротекста, так и IV.

\subsection{Реализация AEAD-схемы}

\subsubsection{Схема Encrypt-then-MAC}

KubeBao использует схему \textbf{Encrypt-then-MAC} (EtM)~--- доказуемо стойкую композицию шифрования и аутентификации при использовании независимых ключей.

\textbf{Теорема (Bellare \& Namprempre, 2000):} Если $\mathrm{Enc}$~--- IND-CPA-стойкая схема шифрования, а $\mathrm{MAC}$~--- SUF-CMA-стойкий код аутентификации, то конструкция EtM с независимыми ключами является IND-CCA-стойкой и INT-CTXT-стойкой.

\subsubsection{Вывод ключей (Key Derivation)}

Из 256-битного мастер-ключа $K_{\mathrm{master}}$ выводятся два независимых 256-битных ключа:
\begin{align}
    K_{\mathrm{enc}} &= \mathrm{SHA\text{-}256}(\text{``kubebao-kuznyechik-enc''} \| \texttt{0x00} \| K_{\mathrm{master}}) \\
    K_{\mathrm{mac}} &= \mathrm{SHA\text{-}256}(\text{``kubebao-kuznyechik-mac''} \| \texttt{0x00} \| K_{\mathrm{master}})
\end{align}

Разделение доменов (domain separation) обеспечивается:
\begin{enumerate}
    \item уникальными строковыми префиксами;
    \item нулевым байтом-разделителем \texttt{0x00}.
\end{enumerate}

Это гарантирует, что $K_{\mathrm{enc}} \neq K_{\mathrm{mac}}$ с подавляющей вероятностью, а компрометация одного ключа не раскрывает другой.

\subsubsection{Формат шифротекста}

\begin{table}[h]
\centering
\caption{Формат выходных данных AEAD}
\begin{tabular}{@{}llp{8cm}@{}}
\toprule
\textbf{Поле} & \textbf{Размер} & \textbf{Описание} \\
\midrule
\texttt{version} & 1 байт & Версия формата (\texttt{0x01}) \\
\texttt{IV} & 16 байт & Случайный инициализационный вектор (CSPRNG) \\
\texttt{ciphertext} & $|P|$ байт & Шифротекст (Кузнечик-CTR) \\
\texttt{CMAC tag} & 16 байт & Тег аутентификации (Кузнечик-CMAC) \\
\bottomrule
\end{tabular}
\end{table}

Overhead: 33 байта ($1 + 16 + 16$).

\subsubsection{Алгоритм шифрования}

\begin{enumerate}
    \item $\mathrm{IV} \leftarrow \mathrm{CSPRNG}(128)$ \hfill\textit{// 16 случайных байт}
    \item $K_{\mathrm{enc}}, K_{\mathrm{mac}} \leftarrow \mathrm{DeriveKeys}(K_{\mathrm{master}})$ \hfill\textit{// SHA-256 с доменным разделением}
    \item $C \leftarrow \text{Кузнечик-CTR}(K_{\mathrm{enc}}, \mathrm{IV}, P)$ \hfill\textit{// ГОСТ~Р~34.13-2015}
    \item $T \leftarrow \text{Кузнечик-CMAC}(K_{\mathrm{mac}}, \mathrm{IV} \| C)$ \hfill\textit{// ГОСТ~Р~34.13-2015}
    \item \textbf{return} $\texttt{0x01} \| \mathrm{IV} \| C \| T$
\end{enumerate}

\subsubsection{Алгоритм дешифрования}

\begin{enumerate}
    \item Разобрать $\mathrm{data} \to \mathrm{version}, \mathrm{IV}, C, T$
    \item Проверить $\mathrm{version} = \texttt{0x01}$
    \item $K_{\mathrm{enc}}, K_{\mathrm{mac}} \leftarrow \mathrm{DeriveKeys}(K_{\mathrm{master}})$
    \item $T' \leftarrow \text{Кузнечик-CMAC}(K_{\mathrm{mac}}, \mathrm{IV} \| C)$
    \item \textbf{Если} $\mathrm{ConstantTimeCompare}(T, T') \neq 1$ \textbf{то} ОШИБКА
    \item $P \leftarrow \text{Кузнечик-CTR}(K_{\mathrm{enc}}, \mathrm{IV}, C)$
    \item \textbf{return} $P$
\end{enumerate}

Шаг~5 использует \textbf{постоянно-временное сравнение} (\texttt{crypto/subtle.ConstantTimeCompare}) для предотвращения timing-атак.

\subsubsection{Свойства безопасности}

\begin{table}[h]
\centering
\caption{Гарантии безопасности AEAD-схемы}
\begin{tabular}{@{}lp{10cm}@{}}
\toprule
\textbf{Свойство} & \textbf{Гарантия} \\
\midrule
Конфиденциальность (IND-CPA) & Кузнечик-CTR с уникальным IV \\
Целостность (INT-CTXT) & $\mathrm{CMAC}(\mathrm{IV} \| C)$ покрывает все изменяемые поля \\
Аутентичность (EUF-CMA) & 128-битный CMAC-тег \\
Стойкость к подмене IV & IV включён во вход CMAC \\
Стойкость к timing-атакам & \texttt{ConstantTimeCompare} при проверке тега \\
Независимость ключей & Доменное разделение при выводе \\
\bottomrule
\end{tabular}
\end{table}

\subsection{Управление ключами}

\subsubsection{Хранение в OpenBao KV}

Ключи хранятся в OpenBao KV v2 по пути:
\begin{verbatim}
secret/data/{kvPathPrefix}/{keyName}
\end{verbatim}

Формат записи:
\begin{lstlisting}[language=]
{
  "key": "<base64-encoded-256-bit-key>",
  "version": 1
}
\end{lstlisting}

\subsubsection{Кэширование и защита в памяти}

\texttt{KeyManager} реализует:
\begin{enumerate}
    \item \textbf{Двойную проверку блокировки} (double-checked locking) при чтении ключа: \texttt{RLock} $\to$ проверка кэша $\to$ \texttt{RUnlock}; при промахе: \texttt{Lock} $\to$ повторная проверка $\to$ чтение из OpenBao $\to$ \texttt{Unlock}.
    \item \textbf{Копирование ключа} при выдаче: вызывающий код получает копию, а не ссылку на кэш.
    \item \textbf{Зануление памяти} (\texttt{zeroSlice}, \texttt{zeroBytes}) после использования ключевого материала через \texttt{defer}.
    \item \textbf{Инвалидация кэша} (\texttt{InvalidateCache})~--- зануляет ключ в памяти перед удалением ссылки.
\end{enumerate}

\subsubsection{Ротация ключей}

При записи нового ключа в OpenBao KV (с \texttt{version: N+1}):
\begin{enumerate}
    \item \texttt{healthCheckLoop} (каждые 30 секунд) обнаруживает изменение версии;
    \item обновляется \texttt{keyID} в KMS-сервере;
    \item Kubernetes вызывает re-encryption секретов при изменении \texttt{keyID} в \texttt{StatusResponse};
    \item новые секреты шифруются новым ключом;
    \item старый ключ остаётся доступным для дешифрования существующих секретов.
\end{enumerate}

\subsection{Криптографическая стойкость}

\begin{table}[h]
\centering
\caption{Параметры криптографической стойкости}
\begin{tabular}{@{}llp{7cm}@{}}
\toprule
\textbf{Параметр} & \textbf{Значение} & \textbf{Обоснование} \\
\midrule
Стойкость ключа & 256 бит & Квантово-устойчивый ($\geq 128$ бит post-quantum) \\
Стойкость блока & 128 бит & MDS-свойство $L$: 17 активных S-box \\
Стойкость CMAC & 128 бит & Birthday bound: $2^{64}$ запросов \\
Стойкость CTR & $2^{64}$ блоков/IV & Один IV $\to$ $\approx 295$ ЭБ \\
Вывод ключей & SHA-256 & Доменное разделение; PRF-свойство \\
\bottomrule
\end{tabular}
\end{table}

\end{document}
